\chapter{Measuring Performance Like a Scientist}
\label{app:perf}

A game that ``feels fast'' is an opinion; a game that spends 4.2 milliseconds per
frame in \code{update} is a measurement. This appendix shows how to turn the first
into the second using \prg's optional performance-metrics pipeline, which records
per-frame timing and uploads it to a small local server for analysis
\citep{prg32metrics,prg32scimeasure}. The workflow is designed to be reproducible,
so that a classmate or a reviewer can repeat your experiment and obtain the same
result---the defining property of a scientific measurement.

\section{Why measure, and what stays unchanged}

Metrics are disabled by default, so ordinary gameplay is never slowed by
measurement \citep{prg32metrics}. You enable them only for a profiling lab or a
report. When enabled, the runtime begins a measurement \emph{run} when a cartridge
starts and ends it when the cartridge is replaced, recording a sample every few
frames into an in-memory ring buffer that a background task uploads. Because the
recording function merely copies a sample and returns, the game's own timing is
disturbed as little as possible \citep{prg32metrics}.

\section{What each frame records}

A sample captures where a frame's time went and how healthy the system was
\citep{prg32metrics}. The most useful fields for a first report are the three
phase timings and the deadline flag.

\begin{center}
\begin{tabularx}{\linewidth}{@{}lX@{}}
\toprule
Field & Meaning \\
\midrule
\code{frame} & the frame counter \\
\code{update\_us} & microseconds spent in your \code{update} \\
\code{draw\_us} & microseconds spent in your \code{draw} \\
\code{present\_us} & microseconds spent displaying the frame \\
\code{frame\_us} & total work: update + draw + present \\
\code{fps\_x100} & active-work frames per second, times 100 \\
\code{heap\_free} & free memory at that moment \\
\code{deadline\_missed} & true when the frame's work exceeded the 30-FPS budget \\
\bottomrule
\end{tabularx}
\captionof{table}{Per-frame sample fields \citep{prg32metrics}.}
\end{center}

\begin{prgconcept}
The frame budget at 30 frames per second is about 33 milliseconds. If
\code{update}, \code{draw}, and \code{present} together exceed that,
\code{deadline\_missed} is set and the game cannot keep its rate. Profiling is the
art of finding which of the three phases is the culprit and reducing it---a habit
worth forming early.
\end{prgconcept}

\section{The reproducible workflow}

The platform's scientific-measurement tutorial prescribes a discipline borrowed
from experimental science: state the question, freeze the software, measure, and
report enough that someone else can repeat it \citep{prg32scimeasure}.

\paragraph{State a narrow question.} Good questions are observable and repeatable:
how much frame time does each demo spend in update, draw, and present? Does QEMU
produce the same visual output as the board? Which examples stay below the 30-FPS
budget? Write the question, one expected result, and one reason it might be wrong,
\emph{before} measuring \citep{prg32scimeasure}.

\paragraph{Freeze the software state.} Record the exact commit, the ESP-IDF
version, the target (ESP32\nobreakdash-C6 hardware or ESP32\nobreakdash-C3 QEMU), the display backend,
the metrics configuration, and the cartridge build command, so the experiment can
be reconstructed \citep{prg32scimeasure}.

\begin{lstlisting}[style=bash,caption={Recording the software state for reproducibility.}]
git rev-parse HEAD          # the exact source revision
git diff > experiment.patch # any local changes, saved as a patch
\end{lstlisting}

\paragraph{Build with metrics enabled.} A checked-in configuration profile turns
on metrics without disturbing the normal board defaults; you combine the two
profiles in one build \citep{prg32metrics}.

\begin{lstlisting}[style=bash,caption={Building the board firmware with metrics enabled, using both configuration profiles.}]
idf.py -B build-esp32c6-metrics \
  -D SDKCONFIG_DEFAULTS="sdkconfig.defaults;sdkconfig.defaults.metrics" \
  set-target esp32c6
idf.py -B build-esp32c6-metrics \
  -D SDKCONFIG_DEFAULTS="sdkconfig.defaults;sdkconfig.defaults.metrics" \
  build
\end{lstlisting}

\paragraph{Run the metrics server.} A small Flask and SQLite server collects the
uploaded samples \citep{prg32metrics}.

\begin{lstlisting}[style=bash,caption={Starting the local metrics server.}]
python3 -m venv .venv
. .venv/bin/activate
python3 -m pip install -r tools/prg32_metrics_server/requirements.txt
python3 tools/prg32_metrics_server/app.py --host 0.0.0.0 --port 8080
\end{lstlisting}

Point \code{CONFIG\_PRG32\_METRICS\_SERVER\_URL} at the computer's address that the
board or QEMU can reach, then run your cartridge; the firmware uploads samples in
the background as it plays \citep{prg32metrics}.

\section{Getting results out}

The server stores each run and computes summary statistics. You can list runs,
inspect one with its statistics, download the raw samples as CSV, or download a
ready-made Markdown report \citep{prg32metrics}.

\begin{center}
\begin{tabularx}{\linewidth}{@{}lX@{}}
\toprule
Endpoint & Returns \\
\midrule
\code{GET /api/runs} & all recorded runs \\
\code{GET /api/runs/<id>} & one run with summary statistics \\
\code{GET /api/runs/<id>/samples.csv} & the raw per-frame samples \\
\code{GET /api/runs/<id>/report.md} & a formatted report \\
\bottomrule
\end{tabularx}
\captionof{table}{Metrics-server endpoints for retrieving results
\citep{prg32metrics}.}
\end{center}

For a written report there is an export helper that produces, from one run, a CSV
file, a Markdown report, a small \LaTeX{} table you can paste into a document, and,
when \code{matplotlib} is installed, PNG plots \citep{prg32metrics}.

\begin{lstlisting}[style=bash,caption={Exporting a run into report-ready artefacts.}]
python3 tools/prg32_metrics_server/export_run.py <run_id> \
  --db tools/prg32_metrics_server/metrics.db \
  --out metrics_export/<run_id>
\end{lstlisting}

\section{A complete profiling exercise}

The platform suggests a self-contained experiment that teaches both profiling and
the effect of measurement itself on the thing measured \citep{prg32metrics}. Run a
cartridge with metrics off and note the frame rate. Enable metrics with a sample
every frame and run the same cartridge for thirty seconds. Export the run and
compare the update, draw, and present times. Then raise the sample period to one in
five frames and repeat. Finally, explain how the measurement overhead and the
network conditions changed the numbers.

\begin{prgtargets}
\textbf{Both targets.} The metrics pipeline works for both the ESP32\nobreakdash-C6 and the
QEMU target, and the screenshot endpoint returns the same 320-by-240 BMP in either
case, which makes ``does QEMU match the board?'' a question you can answer with
evidence rather than impression \citep{prg32qemu,prg32scimeasure}.
\end{prgtargets}

\begin{prgaside}
The instruction to record one expected result, and one reason it might be wrong,
before measuring is not bureaucracy---it is the heart of honest experiment. A
prediction made after seeing the data is no prediction at all. By committing to an
expectation first, you turn a casual look at some numbers into a genuine test, and
you learn the most from the cases where the machine surprises you.
\end{prgaside}
